% Compile with: uplatex file.tex && dvipdfmx file.dvi
%           or: latexmk -uplatex file.tex
%\documentclass[uplatex,dvipdfmx]{jsarticle}
\documentclass[dvipdfmx,a4paper,11pt]{jsarticle}
\usepackage{amsmath, amssymb}
\usepackage{tikz}

\title{scenario\_line\_balancing}
\date{\today}

\begin{document}
\maketitle

全作業要素の所要時間の総和 $T$ とサイクルタイム（1工程あたりの許容時間）$c$ から、必要な最小工程数 $n=\lceil T/c\rceil$（天井関数、$T/c$ 未満では作れないため切り上げ）を定義する。

$n = ceiling(T / c)$

\begin{equation}
  n = ceiling(T / c)
  \label{eq:E1}
\end{equation}

ライン効率 $\eta=T/(nc)$（理想的な総稼働時間 $T$ に対する実際の総稼働可能時間 $nc$ の比）を定義する。

$\eta = T / (n c)$

\begin{equation}
  \eta = T / (n c)
  \label{eq:E2}
\end{equation}

具体例として総作業時間 $T=88$ を工程数の式に代入する。

E1 の T に 88 を代入して

\begin{equation}
  n = \left \lceil \frac{88}{c} \right \rceil
  \label{eq:E3}
\end{equation}

サイクルタイム $c=15$ も代入し、最小工程数 $n$ を数値として確定する。

E3 の c に 15 を代入して

\begin{equation}
  n = 6
  \label{eq:E4}
\end{equation}

同じ総作業時間 $T=88$ を、今度は効率の式へ代入する。

E2 の T に 88 を代入して

\begin{equation}
  \eta = \frac{88}{\left \lceil \frac{88}{c} \right \rceil c}
  \label{eq:E5}
\end{equation}

同じサイクルタイム $c=15$ を代入し、ライン効率 $\eta$ を数値として確定する。

E5 の c に 15 を代入して

\begin{equation}
  \eta = \frac{44}{45}
  \label{eq:E6}
\end{equation}

バランス遅れ（アイドル率）$d=1-\eta$ を、ステップ6で確定した $\eta$ の値を直接埋め込んで計算する。

$d = 1 - E6$

\begin{equation}
  d = \frac{1}{45}
  \label{eq:E7}
\end{equation}

\end{document}